% ------------------------------------------------------------------------------
% - Package Loading
% ------------------------------------------------------------------------------

\usepackage{amsmath}
\usepackage{amsfonts}
\usepackage{amssymb}
\usepackage{amsthm}
\usepackage{framed} % for putting a frame around the ruleset
\usepackage{thmtools} % see: https://tex.stackexchange.com/a/738948
\usepackage{subcaption}

\usepackage{tikz}
\usetikzlibrary{quantikz2}

\usepackage{algorithm} % for pseudocode
\usepackage[noend]{algpseudocode} % for pseudocode
\renewcommand{\algorithmicrequire}{\textbf{Input:}} % redefine require keyword as 'Input:'
\renewcommand{\algorithmicensure}{\textbf{Output:}} % redefine ensure keyword as 'Output:'

\usepackage{xfrac} % nice fractions with \sfrac{}{} in quantikz
\usepackage{lmodern} % fixes special fonts in some diagrams
\usepackage{anyfontsize} % fixes small fonts in some diagrams
\usepackage{braket} % braket notation
\usepackage{xspace} % spaces after relation names

\usepackage{xcolor} % highlighting and editing

\usepackage{hyperref}
\usepackage[capitalise]{cleveref}

\DeclareMathOperator{\spn}{span}

% ------------------------------------------------------------------------------
% - mathabx: bigboxtimes
% ------------------------------------------------------------------------------

\DeclareFontFamily{U}{mathx}{\hyphenchar\font45}
\DeclareFontShape{U}{mathx}{m}{n}{
      <5> <6> <7> <8> <9> <10>
      <10.95> <12> <14.4> <17.28> <20.74> <24.88>
      mathx10
      }{}
\DeclareSymbolFont{mathx}{U}{mathx}{m}{n}
\DeclareFontSubstitution{U}{mathx}{m}{n}

\DeclareMathSymbol{\bigboxtimes}{1}{mathx}{"D2}

% ------------------------------------------------------------------------------
% - Relation Name Macros
% ------------------------------------------------------------------------------

\newcommand{\HqcHDef}{\text{E(1)}\xspace}
\newcommand{\HqcTargDef}{\text{E(2)}\xspace}
\newcommand{\HqcOPhaseDef}{\text{E(3)}\xspace}
\newcommand{\HqcTargMDef}{\text{E(4)}\xspace}
\newcommand{\HqcSwapDef}{\text{E(5)}\xspace}
\newcommand{\HqcEuler}{\text{E(6)}\xspace}
\newcommand{\HqcPOrd}{\text{E(7)}\xspace}
\newcommand{\HqcCtrlId}{\text{E(8)}\xspace}
\newcommand{\HqcCtrlComp}{\text{E(9)}\xspace}
\newcommand{\HqcCtrlSwap}{\text{E(10)}\xspace}
\newcommand{\HqcFreePhase}{\text{E(11)}\xspace}
\newcommand{\HqcCtrlDag}{\text{E(12)}\xspace}
\newcommand{\HqcCtrlPair}{\text{E(13)}\xspace}
\newcommand{\HqcCtrlExpand}{\text{E(14)}\xspace}
\newcommand{\HqcCtrlMerge}{\text{E(15)}\xspace}
\newcommand{\HqcPowZero}{\text{E(16)}\xspace}
\newcommand{\HqcPowId}{\text{E(17)}\xspace}
\newcommand{\HqcPowAdd}{\text{E(18)}\xspace}
\newcommand{\HqcPowIdem}{\text{E(19)}\xspace}
\newcommand{\HqcPowDag}{\text{E(20)}\xspace}
\newcommand{\HqcPowExpand}{\text{E(21)}\xspace}
\newcommand{\HqcPowCong}{\text{E(22)}\xspace}
\newcommand{\HQCCtrlCong}{\text{E(23)}\xspace}

% ------------------------------------------------------------------------------
% - Theorem Formatting
% ------------------------------------------------------------------------------

\theoremstyle{plain}
\newtheorem{theorem}{Theorem}[section]
\newtheorem{lemma}[theorem]{Lemma}
\newtheorem{proposition}[theorem]{Proposition}
\newtheorem{corollary}[theorem]{Corollary}

\theoremstyle{definition}
\newtheorem{definition}[theorem]{Definition}
\newtheorem{example}[theorem]{Example}
\newtheorem{remark}[theorem]{Remark}

% ------------------------------------------------------------------------------
% - Editing Macros
% ------------------------------------------------------------------------------

\newcommand{\todo}[1]{{\color{red} #1}}
